\chapter{Unit 4: Ideation and Creative Problem Solving}


\section{Section A: Multiple Choice Practice Questions (60 Curated Problems)}

\noindent Practice and test your mastery across all core concepts of this unit. Each question reflects standard university examination patterns and Bloom's Taxonomy competencies.
\vspace{3mm}

\noindent\textbf{Q1. What is the main purpose of exploring creative thinking methods?} \hfill [Tag: Exploration of creative thinking methods for generating solution alternatives | Easy]
\begin{enumerate}[label=(\alph*),topsep=2pt,itemsep=1pt,partopsep=0pt,parsep=0pt]
    \item {To choose the cheapest solution immediately}
    \item {To avoid understanding the problem}
    \item {To copy an existing solution}
    \item {To generate different possible solutions}
\end{enumerate}
\vspace{2mm}

\noindent\textbf{Q2. Which approach encourages people to challenge usual assumptions?} \hfill [Tag: Exploration of creative thinking methods for generating solution alternatives | Easy]
\begin{enumerate}[label=(\alph*),topsep=2pt,itemsep=1pt,partopsep=0pt,parsep=0pt]
    \item {Routine planning}
    \item {Lateral thinking}
    \item {Repeated testing}
    \item {Conventional thinking}
\end{enumerate}
\vspace{2mm}

\noindent\textbf{Q3. What does divergent thinking mainly involve?} \hfill [Tag: Exploration of creative thinking methods for generating solution alternatives | Easy]
\begin{enumerate}[label=(\alph*),topsep=2pt,itemsep=1pt,partopsep=0pt,parsep=0pt]
    \item {Rejecting all unusual ideas}
    \item {Generating many possible ideas}
    \item {Ranking completed solutions}
    \item {Finding one final answer}
\end{enumerate}
\vspace{2mm}

\noindent\textbf{Q4. What does convergent thinking mainly involve?} \hfill [Tag: Exploration of creative thinking methods for generating solution alternatives | Easy]
\begin{enumerate}[label=(\alph*),topsep=2pt,itemsep=1pt,partopsep=0pt,parsep=0pt]
    \item {Changing the problem statement}
    \item {Selecting and refining ideas}
    \item {Creating unlimited ideas}
    \item {Avoiding evaluation completely}
\end{enumerate}
\vspace{2mm}

\noindent\textbf{Q5. What is brainstorming?} \hfill [Tag: Techniques for idea generation | Easy]
\begin{enumerate}[label=(\alph*),topsep=2pt,itemsep=1pt,partopsep=0pt,parsep=0pt]
    \item {A method for measuring user satisfaction}
    \item {A method for generating many ideas}
    \item {A method for writing final instructions}
    \item {A method for calculating project costs}
\end{enumerate}
\vspace{2mm}

\noindent\textbf{Q6. What is an important rule during an early brainstorming session?} \hfill [Tag: Techniques for idea generation | Easy]
\begin{enumerate}[label=(\alph*),topsep=2pt,itemsep=1pt,partopsep=0pt,parsep=0pt]
    \item {Criticize each idea immediately}
    \item {Select the best idea before continuing}
    \item {Allow only one idea per person}
    \item {Focus first on generating ideas}
\end{enumerate}
\vspace{2mm}

\noindent\textbf{Q7. What does SCAMPER help people do?} \hfill [Tag: Techniques for idea generation | Easy]
\begin{enumerate}[label=(\alph*),topsep=2pt,itemsep=1pt,partopsep=0pt,parsep=0pt]
    \item {Measure the time of a meeting}
    \item {Replace user research completely}
    \item {Identify the final project budget}
    \item {Improve or change existing ideas}
\end{enumerate}
\vspace{2mm}

\noindent\textbf{Q8. What is mind mapping used for?} \hfill [Tag: Techniques for idea generation | Easy]
\begin{enumerate}[label=(\alph*),topsep=2pt,itemsep=1pt,partopsep=0pt,parsep=0pt]
    \item {Writing a legal agreement}
    \item {Organizing related thoughts visually}
    \item {Testing a physical prototype}
    \item {Calculating a market price}
\end{enumerate}
\vspace{2mm}

\noindent\textbf{Q9. What is the purpose of using a random word in idea generation?} \hfill [Tag: Techniques for idea generation | Easy]
\begin{enumerate}[label=(\alph*),topsep=2pt,itemsep=1pt,partopsep=0pt,parsep=0pt]
    \item {To prove an idea is correct}
    \item {To identify the project owner}
    \item {To introduce a fresh connection}
    \item {To remove all project criteria}
\end{enumerate}
\vspace{2mm}

\noindent\textbf{Q10. What does the phrase "How might we...?" usually introduce?} \hfill [Tag: Techniques for idea generation | Easy]
\begin{enumerate}[label=(\alph*),topsep=2pt,itemsep=1pt,partopsep=0pt,parsep=0pt]
    \item {A list of project expenses}
    \item {A completed design decision}
    \item {A final evaluation score}
    \item {An opportunity for solutions}
\end{enumerate}
\vspace{2mm}

\noindent\textbf{Q11. What is a criterion?} \hfill [Tag: Evaluation of solutions based on defined criteria and objectives | Easy]
\begin{enumerate}[label=(\alph*),topsep=2pt,itemsep=1pt,partopsep=0pt,parsep=0pt]
    \item {A description of the design team}
    \item {A random personal preference}
    \item {A standard used for judging options}
    \item {A list of unrelated ideas}
\end{enumerate}
\vspace{2mm}

\noindent\textbf{Q12. Why should solutions be evaluated against defined objectives?} \hfill [Tag: Evaluation of solutions based on defined criteria and objectives | Easy]
\begin{enumerate}[label=(\alph*),topsep=2pt,itemsep=1pt,partopsep=0pt,parsep=0pt]
    \item {To check whether they meet the goals}
    \item {To avoid comparing alternatives}
    \item {To increase the number of ideas}
    \item {To remove the need for user feedback}
\end{enumerate}
\vspace{2mm}

\noindent\textbf{Q13. Which criterion could be used to evaluate a new product design?} \hfill [Tag: Evaluation of solutions based on defined criteria and objectives | Easy]
\begin{enumerate}[label=(\alph*),topsep=2pt,itemsep=1pt,partopsep=0pt,parsep=0pt]
    \item {Meeting room size}
    \item {Number of emails sent}
    \item {Team member's birthday}
    \item {Ease of use}
\end{enumerate}
\vspace{2mm}

\noindent\textbf{Q14. What is a scoring matrix used for?} \hfill [Tag: Evaluation of solutions based on defined criteria and objectives | Easy]
\begin{enumerate}[label=(\alph*),topsep=2pt,itemsep=1pt,partopsep=0pt,parsep=0pt]
    \item {Recording attendance at meetings}
    \item {Replacing the design process}
    \item {Generating unrelated problems}
    \item {Comparing solutions using criteria}
\end{enumerate}
\vspace{2mm}

\noindent\textbf{Q15. What does feasibility describe?} \hfill [Tag: Evaluation of solutions based on defined criteria and objectives | Easy]
\begin{enumerate}[label=(\alph*),topsep=2pt,itemsep=1pt,partopsep=0pt,parsep=0pt]
    \item {How many words describe an idea}
    \item {How colorful an idea is}
    \item {How practical an idea is}
    \item {How old an idea is}
\end{enumerate}
\vspace{2mm}

\noindent\textbf{Q16. What does prioritisation mean?} \hfill [Tag: Prioritisation of solutions based on defined criteria and objectives | Easy]
\begin{enumerate}[label=(\alph*),topsep=2pt,itemsep=1pt,partopsep=0pt,parsep=0pt]
    \item {Generating options without limits}
    \item {Describing the problem for the first time}
    \item {Ordering options by importance}
    \item {Ignoring all evaluation criteria}
\end{enumerate}
\vspace{2mm}

\noindent\textbf{Q17. Why are solutions prioritised?} \hfill [Tag: Prioritisation of solutions based on defined criteria and objectives | Easy]
\begin{enumerate}[label=(\alph*),topsep=2pt,itemsep=1pt,partopsep=0pt,parsep=0pt]
    \item {To make every option identical}
    \item {To prevent ideas from being compared}
    \item {To avoid setting project goals}
    \item {To focus effort on stronger options}
\end{enumerate}
\vspace{2mm}

\noindent\textbf{Q18. Which factor is commonly considered when prioritising solutions?} \hfill [Tag: Prioritisation of solutions based on defined criteria and objectives | Easy]
\begin{enumerate}[label=(\alph*),topsep=2pt,itemsep=1pt,partopsep=0pt,parsep=0pt]
    \item {Designer handwriting}
    \item {Meeting snack choice}
    \item {Order of team introductions}
    \item {Expected user impact}
\end{enumerate}
\vspace{2mm}

\noindent\textbf{Q19. What is the main purpose of selecting a solution?} \hfill [Tag: Selection of solutions based on defined criteria and objectives | Easy]
\begin{enumerate}[label=(\alph*),topsep=2pt,itemsep=1pt,partopsep=0pt,parsep=0pt]
    \item {To delay decisions indefinitely}
    \item {To generate the first possible idea}
    \item {To choose an option for further development}
    \item {To remove all project objectives}
\end{enumerate}
\vspace{2mm}

\noindent\textbf{Q20. A solution should generally be selected because it:} \hfill [Tag: Selection of solutions based on defined criteria and objectives | Easy]
\begin{enumerate}[label=(\alph*),topsep=2pt,itemsep=1pt,partopsep=0pt,parsep=0pt]
    \item {Best matches the defined objectives}
    \item {Was created by the most senior person}
    \item {Was mentioned first in discussion}
    \item {Has the most unusual name}
\end{enumerate}
\vspace{2mm}

\noindent\textbf{Q21. A team is designing a low-cost water filter and wants to move beyond obvious improvements. Which method would best encourage them to examine the product from unrelated perspectives?} \hfill [Tag: Exploration of creative thinking methods for generating solution alternatives | Medium]
\begin{enumerate}[label=(\alph*),topsep=2pt,itemsep=1pt,partopsep=0pt,parsep=0pt]
    \item {Selecting the cheapest material}
    \item {Benchmarking similar filters}
    \item {Repeating the initial brainstorm}
    \item {Using random-word association}
\end{enumerate}
\vspace{2mm}

\noindent\textbf{Q22. A team uses SCAMPER on a reusable lunch container. Asking, "Can the lid be used as a measuring cup?" mainly applies which SCAMPER action?} \hfill [Tag: Exploration of creative thinking methods for generating solution alternatives | Medium]
\begin{enumerate}[label=(\alph*),topsep=2pt,itemsep=1pt,partopsep=0pt,parsep=0pt]
    \item {Put to another use}
    \item {Adapt}
    \item {Substitute}
    \item {Combine}
\end{enumerate}
\vspace{2mm}

\noindent\textbf{Q23. During a design session, participants first generate ideas independently and then share them with the group. What is the main advantage of this approach?} \hfill [Tag: Exploration of creative thinking methods for generating solution alternatives | Medium]
\begin{enumerate}[label=(\alph*),topsep=2pt,itemsep=1pt,partopsep=0pt,parsep=0pt]
    \item {It ranks ideas automatically}
    \item {It removes the need for evaluation}
    \item {It reduces the effect of groupthink}
    \item {It guarantees technical feasibility}
\end{enumerate}
\vspace{2mm}

\noindent\textbf{Q24. A public library wants to redesign its membership process. Which activity best uses analogy as a creative thinking method?} \hfill [Tag: Exploration of creative thinking methods for generating solution alternatives | Medium]
\begin{enumerate}[label=(\alph*),topsep=2pt,itemsep=1pt,partopsep=0pt,parsep=0pt]
    \item {Listing every current complaint}
    \item {Choosing the most familiar software}
    \item {Studying how airports manage passengers}
    \item {Copying another library's form}
\end{enumerate}
\vspace{2mm}

\noindent\textbf{Q25. A facilitator asks participants to produce as many ideas as possible for ten minutes and postpones all criticism. Which brainstorming principle is being applied?} \hfill [Tag: Techniques for idea generation | Medium]
\begin{enumerate}[label=(\alph*),topsep=2pt,itemsep=1pt,partopsep=0pt,parsep=0pt]
    \item {Separate idea generation from evaluation}
    \item {Require immediate feasibility checks}
    \item {Allow only expert suggestions}
    \item {Prioritise quality over quantity}
\end{enumerate}
\vspace{2mm}

\noindent\textbf{Q26. A team has identified the problem "students struggle to find quiet study spaces." Which reframed question is most likely to generate diverse ideas?} \hfill [Tag: Techniques for idea generation | Medium]
\begin{enumerate}[label=(\alph*),topsep=2pt,itemsep=1pt,partopsep=0pt,parsep=0pt]
    \item {Why are students dissatisfied?}
    \item {How can we add more desks?}
    \item {How might we support focused study anywhere?}
    \item {Which room is currently unused?}
\end{enumerate}
\vspace{2mm}

\noindent\textbf{Q27. A team maps the stages a customer experiences when returning an online purchase and generates ideas at each stage. Which technique is this?} \hfill [Tag: Techniques for idea generation | Medium]
\begin{enumerate}[label=(\alph*),topsep=2pt,itemsep=1pt,partopsep=0pt,parsep=0pt]
    \item {Journey-based ideation}
    \item {Root-cause analysis}
    \item {Weighted scoring}
    \item {Forced ranking}
\end{enumerate}
\vspace{2mm}

\noindent\textbf{Q28. In a brainwriting exercise, each participant writes three ideas before passing the sheet to another person. What is the likely benefit?} \hfill [Tag: Techniques for idea generation | Medium]
\begin{enumerate}[label=(\alph*),topsep=2pt,itemsep=1pt,partopsep=0pt,parsep=0pt]
    \item {It gives quieter participants equal input opportunities}
    \item {It prevents ideas from influencing one another}
    \item {It ensures every idea is technically feasible}
    \item {It eliminates the need for a problem statement}
\end{enumerate}
\vspace{2mm}

\noindent\textbf{Q29. A team uses a morphological chart to improve a bicycle-light system. What does the chart mainly help the team do?} \hfill [Tag: Techniques for idea generation | Medium]
\begin{enumerate}[label=(\alph*),topsep=2pt,itemsep=1pt,partopsep=0pt,parsep=0pt]
    \item {Identify the single root cause}
    \item {Combine alternatives for different functions}
    \item {Compare completed prototypes}
    \item {Measure user satisfaction directly}
\end{enumerate}
\vspace{2mm}

\noindent\textbf{Q30. A team evaluates three service concepts using cost, accessibility, and environmental impact. Why should the criteria be defined before scoring the concepts?} \hfill [Tag: Evaluation of solutions based on defined criteria and objectives | Medium]
\begin{enumerate}[label=(\alph*),topsep=2pt,itemsep=1pt,partopsep=0pt,parsep=0pt]
    \item {To reduce arbitrary and inconsistent judgments}
    \item {To avoid involving users in testing}
    \item {To ensure the newest concept wins}
    \item {To make every concept equally expensive}
\end{enumerate}
\vspace{2mm}

\noindent\textbf{Q31. A concept scores 4 out of 5 for usability, 2 for cost, and 5 for safety. If safety is the most important criterion, which evaluation method is most appropriate?} \hfill [Tag: Evaluation of solutions based on defined criteria and objectives | Medium]
\begin{enumerate}[label=(\alph*),topsep=2pt,itemsep=1pt,partopsep=0pt,parsep=0pt]
    \item {Random selection by the team}
    \item {Weighted decision matrix}
    \item {Ranking by implementation speed}
    \item {Unweighted addition of scores}
\end{enumerate}
\vspace{2mm}

\noindent\textbf{Q32. A prototype receives high usability ratings but fails a mandatory safety requirement. How should the team normally treat this result?} \hfill [Tag: Evaluation of solutions based on defined criteria and objectives | Medium]
\begin{enumerate}[label=(\alph*),topsep=2pt,itemsep=1pt,partopsep=0pt,parsep=0pt]
    \item {Average safety with the other scores}
    \item {Select it because usability is high}
    \item {Ignore safety until final production}
    \item {Reject or redesign it because safety is a constraint}
\end{enumerate}
\vspace{2mm}

\noindent\textbf{Q33. Two concepts have similar total scores, but one has much stronger evidence from user testing. What should the team do next?} \hfill [Tag: Evaluation of solutions based on defined criteria and objectives | Medium]
\begin{enumerate}[label=(\alph*),topsep=2pt,itemsep=1pt,partopsep=0pt,parsep=0pt]
    \item {Discard both concepts as inconclusive}
    \item {Choose the concept with more features}
    \item {Select the cheaper concept automatically}
    \item {Examine the evidence quality before deciding}
\end{enumerate}
\vspace{2mm}

\noindent\textbf{Q34. A team evaluates a transport app only on download numbers, although its objective is to improve reliable access for elderly users. What is the main problem?} \hfill [Tag: Evaluation of solutions based on defined criteria and objectives | Medium]
\begin{enumerate}[label=(\alph*),topsep=2pt,itemsep=1pt,partopsep=0pt,parsep=0pt]
    \item {The app has too many evaluation criteria}
    \item {The users are evaluating the app directly}
    \item {The objective is already completely achieved}
    \item {The measure is not aligned with the objective}
\end{enumerate}
\vspace{2mm}

\noindent\textbf{Q35. A team uses an impact-effort matrix. Concept A has high impact and high effort, while Concept B has moderate impact and low effort. Which concept is usually a good early priority?} \hfill [Tag: Prioritisation of solutions based on defined criteria and objectives | Medium]
\begin{enumerate}[label=(\alph*),topsep=2pt,itemsep=1pt,partopsep=0pt,parsep=0pt]
    \item {Neither concept, because impact must be minimal}
    \item {Both concepts, because matrices cannot compare them}
    \item {Concept A, because effort is irrelevant}
    \item {Concept B, because it offers a practical early gain}
\end{enumerate}
\vspace{2mm}

\noindent\textbf{Q36. A product team has limited resources and ranks ideas by customer value, implementation effort, and strategic fit. What is the purpose of this prioritisation?} \hfill [Tag: Prioritisation of solutions based on defined criteria and objectives | Medium]
\begin{enumerate}[label=(\alph*),topsep=2pt,itemsep=1pt,partopsep=0pt,parsep=0pt]
    \item {To generate additional ideas}
    \item {To remove all uncertainty from development}
    \item {To ensure every idea is implemented}
    \item {To focus resources on the most suitable alternatives}
\end{enumerate}
\vspace{2mm}

\noindent\textbf{Q37. A prioritisation matrix gives customer value a weight of 0.5, feasibility 0.3, and sustainability 0.2. Which statement is correct?} \hfill [Tag: Prioritisation of solutions based on defined criteria and objectives | Medium]
\begin{enumerate}[label=(\alph*),topsep=2pt,itemsep=1pt,partopsep=0pt,parsep=0pt]
    \item {Sustainability has the greatest influence}
    \item {Customer value has the greatest influence}
    \item {Feasibility can be ignored}
    \item {All criteria have equal influence}
\end{enumerate}
\vspace{2mm}

\noindent\textbf{Q38. A team must choose one concept for prototyping. Concept X has the highest score, but it depends on technology that is not yet available. What is the most responsible selection approach?} \hfill [Tag: Selection of solutions based on defined criteria and objectives | Medium]
\begin{enumerate}[label=(\alph*),topsep=2pt,itemsep=1pt,partopsep=0pt,parsep=0pt]
    \item {Discard all concepts with technical uncertainty}
    \item {Select the next-highest concept automatically}
    \item {Assess the risk and validate the technology first}
    \item {Select X without recording the dependency}
\end{enumerate}
\vspace{2mm}

\noindent\textbf{Q39. A selection panel chooses a solution because it matches the organisation's stated objective, even though another option is more popular with panel members. Why is this decision defensible?} \hfill [Tag: Selection of solutions based on defined criteria and objectives | Medium]
\begin{enumerate}[label=(\alph*),topsep=2pt,itemsep=1pt,partopsep=0pt,parsep=0pt]
    \item {Popularity is always irrelevant}
    \item {Panel preferences are always biased}
    \item {Selection should reflect defined objectives}
    \item {Objectives matter only after implementation}
\end{enumerate}
\vspace{2mm}

\noindent\textbf{Q40. After selecting a concept, a team documents the criteria, evidence, trade-offs, and assumptions behind its decision. What is the main benefit?} \hfill [Tag: Selection of solutions based on defined criteria and objectives | Medium]
\begin{enumerate}[label=(\alph*),topsep=2pt,itemsep=1pt,partopsep=0pt,parsep=0pt]
    \item {It replaces the need for prototype testing}
    \item {It improves transparency and supports later review}
    \item {It guarantees the concept will succeed}
    \item {It prevents the concept from being modified}
\end{enumerate}
\vspace{2mm}

\noindent\textbf{Q41. A hospital wants to reduce outpatient waiting time without adding staff or rooms. Which sequence best uses creative thinking methods before selecting concepts?} \hfill [Tag: Exploration of creative thinking methods for generating solution alternatives | Hard]
\begin{enumerate}[label=(\alph*),topsep=2pt,itemsep=1pt,partopsep=0pt,parsep=0pt]
    \item {Interview managers, rank their proposals, and implement the top proposal}
    \item {List available technologies, choose the cheapest, and test adoption}
    \item {Benchmark competitors, copy the fastest workflow, and pilot it}
    \item {Define constraints, reframe the problem, diverge broadly, and converge with criteria}
\end{enumerate}
\vspace{2mm}

\noindent\textbf{Q42. A team repeatedly proposes minor improvements to an existing product. Which intervention most directly reduces fixation on the current solution?} \hfill [Tag: Exploration of creative thinking methods for generating solution alternatives | Hard]
\begin{enumerate}[label=(\alph*),topsep=2pt,itemsep=1pt,partopsep=0pt,parsep=0pt]
    \item {Use analogies from unrelated domains and suspend feasibility judgment initially}
    \item {Ask users to score the existing design against competitor products}
    \item {Limit ideation to technologies already available inside the organization}
    \item {Require each member to defend the current design before ideation}
\end{enumerate}
\vspace{2mm}

\noindent\textbf{Q43. A public-transit team reframes "How might we make buses faster?" as "How might we reduce passengers' perceived journey burden?" What is the primary advantage of this reframing?} \hfill [Tag: Exploration of creative thinking methods for generating solution alternatives | Hard]
\begin{enumerate}[label=(\alph*),topsep=2pt,itemsep=1pt,partopsep=0pt,parsep=0pt]
    \item {It guarantees that the eventual solution will reduce travel time}
    \item {It expands the solution space beyond vehicle speed and route changes}
    \item {It removes the need to understand passengers' actual travel behavior}
    \item {It converts a subjective experience into a fully measurable requirement}
\end{enumerate}
\vspace{2mm}

\noindent\textbf{Q44. During a design sprint, senior engineers dominate discussion and junior members stop contributing. Which structure best preserves diversity without sacrificing technical rigor?} \hfill [Tag: Exploration of creative thinking methods for generating solution alternatives | Hard]
\begin{enumerate}[label=(\alph*),topsep=2pt,itemsep=1pt,partopsep=0pt,parsep=0pt]
    \item {Use silent individual ideation, round-robin sharing, and later technical review}
    \item {Assign the most experienced engineer to combine all ideas immediately}
    \item {Separate junior members permanently so their ideas remain unconstrained}
    \item {Allow open debate first, then invite junior members to endorse proposals}
\end{enumerate}
\vspace{2mm}

\noindent\textbf{Q45. A team applies SCAMPER to a reusable food container. Which prompt most specifically represents the "Put to another use" operation?} \hfill [Tag: Techniques for idea generation | Hard]
\begin{enumerate}[label=(\alph*),topsep=2pt,itemsep=1pt,partopsep=0pt,parsep=0pt]
    \item {Can the container be combined with a temperature sensor?}
    \item {Can the shape be reversed to improve stacking?}
    \item {Can the container be used as a seedling tray?}
    \item {Can the lid be replaced with a compostable material?}
\end{enumerate}
\vspace{2mm}

\noindent\textbf{Q46. A team uses morphological analysis for a portable water purifier. It identifies four independent dimensions with 3, 4, 2, and 5 alternatives. How many nominal combinations exist before compatibility constraints are removed?} \hfill [Tag: Techniques for idea generation | Hard]
\begin{enumerate}[label=(\alph*),topsep=2pt,itemsep=1pt,partopsep=0pt,parsep=0pt]
    \item {24}
    \item {120}
    \item {14}
    \item {240}
\end{enumerate}
\vspace{2mm}

\noindent\textbf{Q47. In a reverse-brainstorming session on reducing data breaches, which prompt is most appropriate?} \hfill [Tag: Techniques for idea generation | Hard]
\begin{enumerate}[label=(\alph*),topsep=2pt,itemsep=1pt,partopsep=0pt,parsep=0pt]
    \item {How could we guarantee that every account becomes compromised?}
    \item {What controls would make secure behavior easier for users?}
    \item {What is the average breach rate among comparable organizations?}
    \item {Which existing security feature should receive the largest budget?}
\end{enumerate}
\vspace{2mm}

\noindent\textbf{Q48. A team has generated 200 ideas, but most are variations of three familiar concepts. Which next step best increases useful variety?} \hfill [Tag: Techniques for idea generation | Hard]
\begin{enumerate}[label=(\alph*),topsep=2pt,itemsep=1pt,partopsep=0pt,parsep=0pt]
    \item {Select the three most frequently mentioned concepts for rapid prototyping}
    \item {Ask the team to produce 200 additional ideas under the same prompt}
    \item {Discard all repeated concepts and retain only technically novel proposals}
    \item {Apply random stimuli and forced combinations across different concept families}
\end{enumerate}
\vspace{2mm}

\noindent\textbf{Q49. When using the six thinking hats for a medication-adherence service, which practice best preserves the method's value?} \hfill [Tag: Techniques for idea generation | Hard]
\begin{enumerate}[label=(\alph*),topsep=2pt,itemsep=1pt,partopsep=0pt,parsep=0pt]
    \item {Require every participant to defend one permanent hat throughout the project}
    \item {Use the critical hat first to eliminate unrealistic ideas before alternatives emerge}
    \item {Use one mode at a time so the group shares the same cognitive focus}
    \item {Let participants choose only the hat matching their usual thinking style}
\end{enumerate}
\vspace{2mm}

\noindent\textbf{Q50. A solution must satisfy a legal requirement and then maximize accessibility and cost performance. Which evaluation model is most appropriate?} \hfill [Tag: Evaluation of solutions based on defined criteria and objectives | Hard]
\begin{enumerate}[label=(\alph*),topsep=2pt,itemsep=1pt,partopsep=0pt,parsep=0pt]
    \item {Average stakeholder preferences before checking whether the requirement is met}
    \item {Use a weighted average for all criteria, including the legal requirement}
    \item {Rank solutions by cost first and remove accessibility from the evaluation}
    \item {Screen for the legal requirement first, then use weighted scoring on the remainder}
\end{enumerate}
\vspace{2mm}

\noindent\textbf{Q51. Two concepts receive weighted scores of 82 and 79. The first depends on an untested technology, while the second uses mature technology. What is the most defensible evaluation response?} \hfill [Tag: Evaluation of solutions based on defined criteria and objectives | Hard]
\begin{enumerate}[label=(\alph*),topsep=2pt,itemsep=1pt,partopsep=0pt,parsep=0pt]
    \item {Select the second because mature technology always dominates higher performance}
    \item {Select the first because the numerical score completely resolves uncertainty}
    \item {Average the technology maturity ratings into the existing score without changing weights}
    \item {Represent technology uncertainty explicitly and test whether the ranking is robust}
\end{enumerate}
\vspace{2mm}

\noindent\textbf{Q52. A scoring matrix uses a 1--5 scale for cost, safety, and desirability, but the team assigns weights of 0.5, 0.3, and 0.4. What is the principal flaw?} \hfill [Tag: Evaluation of solutions based on defined criteria and objectives | Hard]
\begin{enumerate}[label=(\alph*),topsep=2pt,itemsep=1pt,partopsep=0pt,parsep=0pt]
    \item {Safety should never be included in a multicriteria evaluation}
    \item {The weights sum to 1.2, so the weighted scores are not properly normalized}
    \item {A five-point scale cannot compare more than two proposed solutions}
    \item {The criteria are measured on an ordinal scale rather than a ratio scale}
\end{enumerate}
\vspace{2mm}

\noindent\textbf{Q53. A concept improves average service speed but creates severe delays for a small, vulnerable user group. Which evaluation change best exposes this issue?} \hfill [Tag: Evaluation of solutions based on defined criteria and objectives | Hard]
\begin{enumerate}[label=(\alph*),topsep=2pt,itemsep=1pt,partopsep=0pt,parsep=0pt]
    \item {Replace all subgroup measures with the overall mean}
    \item {Add distributional or worst-case criteria alongside the average outcome}
    \item {Increase the speed criterion's weight until the average benefit dominates}
    \item {Evaluate only the largest user segment to avoid statistical noise}
\end{enumerate}
\vspace{2mm}

\noindent\textbf{Q54. A product team must prioritize initiatives using value, effort, risk reduction, and strategic alignment. Why is an unweighted impact--effort grid potentially inadequate?} \hfill [Tag: Prioritisation of solutions based on defined criteria and objectives | Hard]
\begin{enumerate}[label=(\alph*),topsep=2pt,itemsep=1pt,partopsep=0pt,parsep=0pt]
    \item {It prevents teams from generating any additional solution alternatives}
    \item {It requires quantitative data that design teams can never collect}
    \item {It guarantees that high-effort solutions will always be rejected}
    \item {It collapses risk and alignment into two dimensions that may hide important trade-offs}
\end{enumerate}
\vspace{2mm}

\noindent\textbf{Q55. A team prioritizes features by multiplying customer value by confidence and dividing by effort. Which interpretation is most accurate?} \hfill [Tag: Prioritisation of solutions based on defined criteria and objectives | Hard]
\begin{enumerate}[label=(\alph*),topsep=2pt,itemsep=1pt,partopsep=0pt,parsep=0pt]
    \item {It is an objective measure that eliminates stakeholder judgment}
    \item {It can be used without checking whether the criteria overlap}
    \item {It proves that low-effort features always have greater strategic value}
    \item {It is a decision rule whose ranking depends strongly on scale definitions}
\end{enumerate}
\vspace{2mm}

\noindent\textbf{Q56. Three solutions are prioritized using a weighted score, but one solution is worse than another on every criterion and no constraint favors it. What should happen?} \hfill [Tag: Prioritisation of solutions based on defined criteria and objectives | Hard]
\begin{enumerate}[label=(\alph*),topsep=2pt,itemsep=1pt,partopsep=0pt,parsep=0pt]
    \item {Remove it as dominated unless missing information could change the comparison}
    \item {Assign it the highest priority to preserve stakeholder disagreement}
    \item {Promote it because uniform weakness indicates a safer compromise}
    \item {Retain it because a weighted score is always sufficient}
\end{enumerate}
\vspace{2mm}

\noindent\textbf{Q57. A final choice must balance user desirability, technical feasibility, and organizational viability. Which selection approach is strongest when evidence is incomplete?} \hfill [Tag: Selection of solutions based on defined criteria and objectives | Hard]
\begin{enumerate}[label=(\alph*),topsep=2pt,itemsep=1pt,partopsep=0pt,parsep=0pt]
    \item {Choose the cheapest concept to minimize irreversible commitment}
    \item {Let the most senior stakeholder select the concept after reviewing the matrix}
    \item {Use staged selection with prototypes designed to resolve critical uncertainties}
    \item {Choose the concept with the highest desirability rating}
\end{enumerate}
\vspace{2mm}

\noindent\textbf{Q58. A concept scores well overall but fails a safety threshold under a foreseeable operating condition. How should it be treated?} \hfill [Tag: Selection of solutions based on defined criteria and objectives | Hard]
\begin{enumerate}[label=(\alph*),topsep=2pt,itemsep=1pt,partopsep=0pt,parsep=0pt]
    \item {Keep it as the preferred option and document the safety concern later}
    \item {Select it because strong performance elsewhere compensates for the failure}
    \item {Average the safety result with positive scores from other criteria}
    \item {Reject or redesign it unless the safety threshold is formally changed}
\end{enumerate}
\vspace{2mm}

\noindent\textbf{Q59. A selection panel's final decision consistently favors concepts proposed by senior staff, although anonymous scoring shows no performance advantage. Which corrective action is most appropriate?} \hfill [Tag: Selection of solutions based on defined criteria and objectives | Hard]
\begin{enumerate}[label=(\alph*),topsep=2pt,itemsep=1pt,partopsep=0pt,parsep=0pt]
    \item {Remove all qualitative judgments from the selection process}
    \item {Use anonymized evidence and independent scoring before group deliberation}
    \item {Select only concepts generated by the largest stakeholder group}
    \item {Increase the senior staff's criterion weights to reflect their experience}
\end{enumerate}
\vspace{2mm}

\noindent\textbf{Q60. A team must choose between a high-potential concept with uncertain adoption and a moderate concept with strong evidence of acceptance. Which decision rule best reflects responsible selection?} \hfill [Tag: Selection of solutions based on defined criteria and objectives | Hard]
\begin{enumerate}[label=(\alph*),topsep=2pt,itemsep=1pt,partopsep=0pt,parsep=0pt]
    \item {Always choose the accepted concept because evidence outweighs future potential}
    \item {Always choose the higher-potential concept because uncertainty implies innovation}
    \item {Compare expected value, uncertainty, reversibility, and the cost of further learning}
    \item {Choose randomly to prevent evidence from suppressing creative alternatives}
\end{enumerate}
\vspace{2mm}


\newpage
\section{Section B: Comprehensive Subjective Problem with Detailed Analytical Rubric}

\begin{subjectiveq}{4}{1}
\textbf{Case Background:} An e-commerce delivery enterprise is confronted with high last-mile carbon footprints, severe traffic congestion, and a 19\% rate of failed residential first-time deliveries in high-density urban corridors. Management seeks breakthrough solution concepts beyond simply buying more vans.

\textbf{Required Design Deliverables:}
\begin{enumerate}[label=(\alph*)]
    \item Apply the SCAMPER creative ideation methodology to generate at least six innovative system-level concepts.
    \item Run a structured 2$\times$2 Impact vs. Effort Prioritization Matrix to position and categorize four distinct candidate solutions.
    \item Apply a Pugh Decision Matrix against a baseline concept to select the optimal design pathway based on Desirability, Feasibility, and Viability.
\end{enumerate}
\end{subjectiveq}

\begin{solbox}
\textbf{1. SCAMPER Ideation Matrix Deployment:}\\
\begin{itemize}
    \item \textbf{Substitute (S):} Replace heavy diesel delivery vans with electric cargo quadricycles and modular walking courier pods.
    \item \textbf{Combine (C):} Combine residential parcel deliveries with existing public bus transit baggage compartments during off-peak hours.
    \item \textbf{Adapt (A):} Adapt the ATM / convenience store model by creating neighborhood parcel locker hubs inside 24/7 metro stations.
    \item \textbf{Modify/Magnify (M):} Magnify recipient autonomy through dynamic 15-minute delivery time-slot selection via instant messaging.
    \item \textbf{Put to Another Use (P):} Use courier delivery vehicles as mobile environmental air-quality sensing nodes during rounds.
    \item \textbf{Eliminate (E):} Eliminate doorstep re-attempts by defaulting failed drops to localized micro-lockers within 200 meters of the residence.
    \item \textbf{Reverse/Rearrange (R):} Invert the delivery flow: crowdsource neighbors who are already home to accept parcels in return for delivery credits.
\end{itemize}

\textbf{2. 2$\times$2 Impact vs. Effort Prioritization Quadrants:}\\
\begin{itemize}
    \item \textbf{Quick Wins (High Impact, Low Effort):} Dynamic SMS 15-minute dispatch alert allowing real-time drop approval (Drastically cuts failed deliveries immediately).
    \item \textbf{Major Strategic Projects (High Impact, High Effort):} Network of solar-powered smart automated parcel lockers across municipal subway stations.
    \item \textbf{Fill-ins / Low-Hanging Fruit (Low Impact, Low Effort):} Packaging size standardizations and QR-code scannable doorstep drop boxes.
    \item \textbf{Thankless Tasks / Money Pits (Low Impact, High Effort):} Autonomous sidewalk drone rovers (plagued by curb-cut regulations, vandalism, and high capital cost).
\end{itemize}

\textbf{3. Pugh Selection Matrix (Baseline = Current Diesel Doorstep Delivery):}\\
{\small
\begin{tabularx}{\linewidth}{|X|c|c|c|c|} \hline
\textbf{Criteria} & \textbf{Baseline} & \textbf{Lockers} & \textbf{E-Bikes} & \textbf{Neighbor} \\ \hline
Carbon Footprint Reduction & 0 & +1 & +1 & +1 \\ \hline
First-Time Success Rate & 0 & +1 & 0 & +1 \\ \hline
Customer Convenience & 0 & 0 & 0 & -1 \\ \hline
Capital Expenditure & 0 & -1 & -1 & +1 \\ \hline
Regulatory Ease & 0 & +1 & 0 & -1 \\ \hline
\textbf{Total Score ($\Sigma + / - / 0$)} & \textbf{0} & \textbf{+2 Net} & \textbf{0 Net} & \textbf{+1 Net} \\ \hline
\end{tabularx}\par}

\textit{Synthesis:} The Smart Neighborhood Locker network achieves the highest net score (+2), delivering superior reliability and carbon savings while avoiding the interpersonal friction of crowdsourced neighbor models.
\end{solbox}

\begin{examinertip}
\textbf{Criteria for Full 10/10 Marks:}\\
1. \textbf{Rigorous Theoretical Framing (3 Marks):} Explicitly identify and define the formal model or framework applied (e.g., d.school stages, 5-Whys, MCDA, SCAMPER, Blueprinting, Minto SCQA). Avoid vague hand-waving.\\
2. \textbf{Contextual Application (4 Marks):} Directly connect every step to the specific facts, users, and constraints given in the case scenario.\\
3. \textbf{Actionable Deliverables \& Invariants (3 Marks):} Provide precise, formatted tables, matrices, or SMART criteria with clear quantitative bounds.
\end{examinertip}
